\section{Conclusiones}

Este trabajo ha demostrado que es posible extender una herramienta \textit{open-source}
de producción audiovisual en directo ---Voctomix--- con funcionalidades de nivel
profesional sin comprometer su arquitectura modular ni su rendimiento en tiempo real.

Las seis funcionalidades implementadas (rotulación dinámica, stream blanker avanzado,
Audio Follows Video, auto-off de overlays, telemetría con el servicio notario y
contenerización Docker) responden a necesidades reales identificadas en producciones
de eventos técnicos, y han sido validadas satisfactoriamente en los tres escenarios
de despliegue propuestos.

Desde el punto de vista técnico, la elección de GStreamer como motor multimedia ha
resultado especialmente acertada: su modelo de \textit{pipeline} dinámico ha
permitido integrar todas las nuevas funcionalidades como módulos independientes que
se instancian y conectan en función de la configuración, sin necesidad de modificar
el núcleo del sistema original.

La contenerización mediante Docker Compose representa la aportación más transversal
del proyecto: elimina la principal barrera de adopción de Voctomix ---la instalación
manual de dependencias GStreamer--- y abre la puerta a despliegues en infraestructura
\textit{cloud} o en clústeres de orquestación de contenedores.

En definitiva, Voctomix 2.0 constituye una plataforma de producción audiovisual en
directo robusta, reproducible y extensible, adecuada tanto para producciones de
eventos técnicos universitarios como para conferencias de mayor escala.
